% Generated draft: requires mathematical and source review.
\documentclass[11pt]{article}
\usepackage[margin=25mm]{geometry}
\usepackage{fontspec}
\setmainfont{Latin Modern Roman}
\usepackage{amsmath,amssymb,mathtools,unicode-math}
\setmathfont{Latin Modern Math}
\usepackage{enumitem,needspace,xurl,hyperref}
\setlength{\parindent}{0pt}
\setlength{\parskip}{5pt}
\setlength{\emergencystretch}{3em}
\newcommand{\N}{\mathbb{N}}
\newcommand{\Z}{\mathbb{Z}}
\newcommand{\Q}{\mathbb{Q}}
\newcommand{\R}{\mathbb{R}}
\newcommand{\C}{\mathbb{C}}
\newcommand{\F}{\mathbb{F}}
\newcommand{\A}{\mathbb{A}}
\newcommand{\PP}{\mathbb{P}}
\newcommand{\E}{\mathbb{E}}
\newcommand{\eps}{\varepsilon}
\newcommand{\dd}{\,\mathrm{d}}
\begin{document}
\section*{UnsolvedMath ID: 9700025}
\textbf{Source problem code:} catalogue entry
\subsection{Definitions and mathematical statement}
In the source's city-growth model, let \(N_{(1)}(t)\) denote the largest city population at model time \(t\), as described by the supplied summary. Fix parameters \(\alpha=1\) and \(\beta>2\). Conjecture 22 asserts
\[
 N_{(1)}(t)=t(\log t)^{1-2/\beta+o(1)}\qquad\text{almost surely as }t\to\infty.
\]
Precisely, with \(\gamma=1-2/\beta\), the logarithmic correction statement is
\[
 \Pr\left(\lim_{t\to\infty}
 \frac{\log(N_{(1)}(t)/t)}{\log\log t}=\gamma\right)=1,
\]
where the limit follows the model's time index and populations are positive for the relevant large times. Equivalently, with probability one, for every \(\varepsilon>0\), there is a finite random \(t_0\) such that
\[
 t(\log t)^{\gamma-\varepsilon}\leq N_{(1)}(t)
 \leq t(\log t)^{\gamma+\varepsilon}\qquad(t\geq t_0).
\]
The \(o(1)\) is in the logarithmic exponent and may depend on the sample path. No ratio-one asymptotic with a specified multiplicative constant is asserted.

The underlying Aldous--Huang model's probability law, city creation and population dynamics, meaning of \(\alpha,\beta\), and time normalization are not included in the excerpt. These must be supplied to define \(N_{(1)}\) and its almost-sure event. The statement is for each fixed admissible \(\beta>2\); no uniformity or shared probability space over all \(\beta\) is added.
\subsection{Short English statement}
At \(\alpha=1\), the largest city is conjectured to grow almost surely like \(t\) times a logarithmic factor whose exponent tends to \(1-2/\beta\), for fixed \(\beta>2\). The exact city model and time scale are missing.
\subsection{Sources}
\begin{itemize}
\item Local UnsolvedMath registry, permanent ID 9700025, source code catalogue entry; David Aldous; no proposal year supplied; updated 2026-08-06.
\item Aldous, \emph{Open problems}, cities.html and technical notes, Conjecture 22: \url{https://www.stat.berkeley.edu/~aldous/Research/OP/index.html}. Registry access date 2026-07-29; HTML text with linked primary-PDF text. Background reports 23 active pages, with delegated explicit subproblems separated and compound bullets retained. Present status and redistribution terms require review.
\item The supplied summary identifies the Aldous--Huang model and cites arXiv:1209.5120, reporting that the \(\alpha=1\) case remains unsettled. No title or exact later reference is supplied; these status and search claims were not independently verified here.
\end{itemize}
\end{document}
